Un elemento bimetálico está formado por dos tiras de metales diferentes (acero con $\alpha_a = 11\times 10^{-6}$$^o$C$^{-1}$  y cinc con $\alpha_c = 25\times 10^{-6}$$^o$C$^{-1}$), firmemente unidas. Suponga que dicha lámina se calienta. Trate de describir qué sucederá al bimetálico en virtud de la dilatación de ambos metales. Haga un dibujo que muestre el aspecto de la lámina después del calentamiento (este dispositivo suele emplearse para cerrar un circuito eléctrico; por ejemplo, en las alarmas de incendio).